\heading{数据结构与算法B-21春-期末}
\noteinfo{题目来源于真题扫描件转录，参考解答由 AI 生成。本次整理提供 C 语言版与 Python 语言版；除程序代码语言外，两版题目内容保持一致。}

\textbf{一、选择题（30 分，每小题 2 分）}

\begin{question}
下列不影响算法时间复杂性的因素有（\quad）。
\begin{enumerate}
  \item 问题的规模
  \item 输入值
  \item 计算结果
  \item 算法的策略
\end{enumerate}
\begin{solution}
选 C。算法时间复杂度由问题规模、输入数据和算法策略决定，与计算结果无关。
\end{solution}
\end{question}

\begin{question}
链表不具有的特点是（\quad）。
\begin{enumerate}
  \item 可随机访问任意元素
  \item 插入和删除不需要移动元素
  \item 不必事先估计存储空间
  \item 所需空间与线性表长度成正比
\end{enumerate}
\begin{solution}
选 A。链表是顺序访问结构，不支持随机访问。
\end{solution}
\end{question}

\begin{question}
设有三个元素 $X,Y,Z$ 顺序进栈（进的过程中允许出栈），下列得不到的出栈排列是（\quad）。
\begin{enumerate}
  \item $XYZ$
  \item $YZX$
  \item $ZXY$
  \item $ZYX$
\end{enumerate}
\begin{solution}
选 C。$ZXY$ 不可能：若 $Z$ 第一个出栈，则 $X,Y$ 已在栈中且 $Y$ 在栈顶，之后必须先出 $Y$ 再出 $X$，不能得到 $ZXY$。
\end{solution}
\end{question}

\begin{question}
判定一个队列 $Q$（申请数组空间为 $m$）为空的条件是（\quad）。
\begin{enumerate}
  \item \texttt{Q.rear == Q.front}
  \item \texttt{Q.front != Q.rear}
  \item \texttt{Q.front == (Q.rear+1)\% m}
  \item \texttt{Q.front != (Q.rear+1)\%m}
\end{enumerate}
\begin{solution}
选 A。循环队列判空条件为 \texttt{rear == front}，判满条件为 \texttt{(rear+1)\%m == front}。
\end{solution}
\end{question}

\begin{question}
若定义二叉树中根结点的层数为 $0$，树的高度等于其结点的最大层数加一。则当某二叉树的前序遍历序列和后序遍历序列正好相反，则该二叉树一定是（\quad）。
\begin{enumerate}
  \item 空或只有一个结点
  \item 高度等于其节点数
  \item 任一结点无左孩子
  \item 任一结点无右孩子
\end{enumerate}
\begin{solution}
选 B。先序和后序相反意味着每层只有一个结点，树退化为单链，高度等于节点数。
\end{solution}
\end{question}

\begin{question}
任意一棵二叉树中，所有叶结点在前序遍历序列、中序遍历序列和后序遍历序列中的相对次序（\quad）。
\begin{enumerate}
  \item 发生改变
  \item 不发生改变
  \item 不能确定
  \item 以上都不对
\end{enumerate}
\begin{solution}
选 B。叶子节点在三种遍历中均保持从左到右的相对顺序。
\end{solution}
\end{question}

\begin{question}
假设线性表中每个元素有两个数据项 key1 和 key2，现对线性表按以下规则进行排序：先根据 key1 的值进行非递减排序；在 key1 值相同的情况下，再根据 key2 的值进行非递减排序。满足这种要求的排序方法是（\quad）。
\begin{enumerate}
  \item 先按 key1 冒泡排序，再按 key2 直接选择排序
  \item 先按 key2 冒泡排序，再按 key1 直接选择排序
  \item 先按 key1 直接选择排序，再按 key2 冒泡排序
  \item 先按 key2 直接选择排序，再按 key1 冒泡排序
\end{enumerate}
\begin{solution}
选 B。需要稳定排序：先按次要关键字（key2）排序，再按主要关键字（key1）稳定排序，保证 key1 相同时 key2 有序。
\end{solution}
\end{question}

\begin{question}
有 $n$ 个整数，找到其中最小整数需要比较次数至少是（\quad）次。
\begin{enumerate}
  \item $n$
  \item $\log_2 n$
  \item $n^2-1$
  \item $n-1$
\end{enumerate}
\begin{solution}
选 D。找最小值至少需要 $n-1$ 次比较（每个元素至少参与一次比较，锦标赛形式）。
\end{solution}
\end{question}

\begin{question}
$n$ 个顶点的无向完全图的边数为（\quad）。
\begin{enumerate}
  \item $n(n-1)$
  \item $n(n+1)$
  \item $n(n-1)/2$
  \item $n(n+1)/2$
\end{enumerate}
\begin{solution}
选 C。无向完全图每对顶点间有一条边，共 $C_n^2 = n(n-1)/2$ 条。
\end{solution}
\end{question}

\begin{question}
设无向图 $G=(V,E)$，$G'=(V',E')$，如果 $G'$ 是 $G$ 的生成树，则下面说法中错误的是（\quad）。
\begin{enumerate}
  \item $G'$ 是 $G$ 的子图
  \item $G'$ 是 $G$ 的连通分量
  \item $G'$ 是 $G$ 的极小连通子图且 $V=V'$
  \item $G'$ 是 $G$ 的一个无环子图
\end{enumerate}
\begin{solution}
选 B。生成树是包含所有顶点的极小连通子图，但不一定是连通分量（连通分量是极大连通子图）。
\end{solution}
\end{question}

\begin{question}
有一个散列表，其散列函数为 $h(\text{key})=\text{key}\bmod 13$，使用再散列函数 $h_2(\text{key})=\text{key}\bmod 3$ 解决碰撞。从表中检索出关键码 $38$ 需进行几次比较（\quad）。
\begin{enumerate}
  \item $1$
  \item $2$
  \item $3$
  \item $4$
\end{enumerate}
\begin{solution}
选 B。$h(38)=38\bmod13=12$，冲突后 $h_2(38)=38\bmod3=2$，探查 $(12+2)\%13=1$，找到。
\end{solution}
\end{question}

\begin{question}
在一棵度为 $3$ 的树中，度为 $3$ 的节点个数为 $2$，度为 $2$ 的节点个数为 $1$，则度为 $0$ 的节点个数为（\quad）。
\begin{enumerate}
  \item $4$
  \item $5$
  \item $6$
  \item $7$
\end{enumerate}
\begin{solution}
选 C。树中总边数 $=$ 总度数 $=3\times2+2\times1+n_1$，总节点数 $=2+1+n_1+n_0$，由边数 $=$ 节点数 $-1$ 得 $n_0=6$。
\end{solution}
\end{question}

\begin{question}
由同一组关键字集合构造的各棵二叉排序树（\quad）。
\begin{enumerate}
  \item 其形态不一定相同，但平均查找长度相同
  \item 其形态不一定相同，平均查找长度也不一定相同
  \item 其形态均相同，但平均查找长度不一定相同
  \item 其形态均相同，平均查找长度也都相同
\end{enumerate}
\begin{solution}
选 B。BST 形态取决于插入顺序，不同形态影响查找效率。
\end{solution}
\end{question}

\begin{question}
允许表达式内多种括号混合嵌套，检查表达式中括号是否正确配对的算法，通常选用（\quad）。
\begin{enumerate}
  \item 栈
  \item 线性表
  \item 队列
  \item 二叉排序树
\end{enumerate}
\begin{solution}
选 A。栈的后进先出特性适合括号匹配检查。
\end{solution}
\end{question}

\begin{question}
下列选项中最适合实现数据的频繁增删及高效查找的组织结构是（\quad）。
\begin{enumerate}
  \item 有序表
  \item 堆排序
  \item 二叉排序树
  \item 快速排序
\end{enumerate}
\begin{solution}
选 C。BST 支持 $O(\log n)$ 的查找与 $O(\log n)$ 的插入/删除（平均情况）。
\end{solution}
\end{question}

\clearpage

\textbf{二、判断题（10 分，每小题 1 分；对 Y，错 N）}

\begin{question}
考虑一个长度为 $n$ 的顺序表中各个位置插入新元素的概率相同，则顺序表的插入算法平均时间复杂度为 $O(n)$。（\quad）
\begin{solution}
Y。平均需移动 $n/2$ 个元素。
\end{solution}
\end{question}

\begin{question}
希尔排序算法的每一趟都要调用一次或多次直接插入排序算法，所以其效率比直接插入排序算法差。（\quad）
\begin{solution}
N。希尔排序通过分组使数组逐步有序，总体效率优于直接插入排序。
\end{solution}
\end{question}

\begin{question}
直接插入排序、冒泡排序、Shell 排序都是在数据正序的情况下比数据在逆序的情况下要快。（\quad）
\begin{solution}
Y。正序时比较和移动次数最少。
\end{solution}
\end{question}

\begin{question}
由于碰撞的发生，基于散列表的检索仍然需要进行关键码对比，并且关键码的比较次数仅取决于选择的散列函数与处理碰撞的方法两个因素。（\quad）
\begin{solution}
N。还与装填因子（表满程度）有关。
\end{solution}
\end{question}

\begin{question}
若有一个叶子结点是二叉树中某个子树的前序遍历结果序列的最后一个结点，则它一定是该子树的中序遍历结果序列的最后一个结点。（\quad）
\begin{solution}
Y。前序最后一个节点是某子树的最右叶节点，中序中它也是该子树最后一个节点。
\end{solution}
\end{question}

\begin{question}
若某非空二叉树的前序遍历序列和后序遍历序列正好相同，则该二叉树只有一个根结点。（\quad）
\begin{solution}
Y。先序和后序相同说明每个节点都没有左右子树，即只有一个根节点。
\end{solution}
\end{question}

\begin{question}
有 $n$ 个节点的二叉排序树有多种，其中树高最小的二叉排序树是搜索效率最好。（\quad）
\begin{solution}
Y。平衡的 BST（树高 $O(\log n)$）搜索效率最高。
\end{solution}
\end{question}

\begin{question}
强连通分量是有向图中的最小强连通子图。（\quad）
\begin{solution}
N。强连通分量是极大强连通子图，不是最小。
\end{solution}
\end{question}

\begin{question}
用邻接矩阵法存储一个图时，在不考虑压缩存储的情况下，所占用的存储空间大小只与图中结点个数有关，而与图的边数无关。（\quad）
\begin{solution}
Y。邻接矩阵大小为 $n\times n$，与边数无关。
\end{solution}
\end{question}

\begin{question}
若定义一个有向图的根是指可以从这个结点到达图中任意其他结点，则可知一个有向图中至少有一个根。（\quad）
\begin{solution}
N。有向图可能没有根（如存在多个强连通分量且无法互达）。
\end{solution}
\end{question}

\clearpage

\textbf{三、填空题（20 分，每题 2 分）}

\begin{question}
线性表的顺序存储与链式存储是两种常见存储形式；当表元素有序排序进行二分检索时，应采用 \underline{\hspace{4em}} 存储形式。
\begin{solution}
二分查找需要随机访问中间元素，顺序表支持 $O(1)$ 随机访问，链表不支持。
\end{solution}
\end{question}

\begin{question}
设环形队列的容量为 $20$（单元编号 $0$ 到 $19$），经过一系列的入队和出队运算后，队头变量 front=18，队尾变量 rear=11，则环形队列中有 \underline{\hspace{4em}} 个元素。
\begin{solution}
元素数 $= (rear - front + 20) \% 20 = (11 - 18 + 20) \% 20 = 13$。
\end{solution}
\end{question}

\begin{question}
一棵含有 $101$ 个结点的二叉树中有 $36$ 个叶子结点，度为 $2$ 的结点个数是 \underline{\hspace{4em}}，度为 $1$ 的结点个数是 \underline{\hspace{4em}}。
\begin{solution}
$n_0=36$，由二叉树性质 $n_0=n_2+1$ 得 $n_2=35$，$n_1=101-36-35=30$。
\end{solution}
\end{question}

\begin{question}
已知二叉树的前序遍历结果为 $ADC$，这棵二叉树的树型有 \underline{\hspace{4em}} 种可能。
\begin{solution}
不同树型数目为 Catalan 数 $C_3=5$。
\end{solution}
\end{question}

\begin{question}
已知二叉树的中序序列为 $DGBAECF$，后序序列为 $GDBEFCA$，该二叉树的先序序列是 \underline{\hspace{4em}}。
\begin{solution}
由后序确定根 $A$，中序划分左右子树，递归可得先序：$A,B,D,G,C,E,F$。
\end{solution}
\end{question}

\begin{question}
对于具有 $57$ 个结点的完全二叉树，按层次自顶向下、同一层自左向右、从 $0$ 开始编号。编号为 $18$ 的结点的父结点编号为 \underline{\hspace{4em}}，编号为 $19$ 的结点的右子女结点编号为 \underline{\hspace{4em}}。
\begin{solution}
完全二叉树节点 $i$ 的父节点为 $\lfloor(i-1)/2\rfloor$，左子为 $2i+1$。$18$ 的父：$(18-1)/2=8$；$19$ 的右子：$2\times19+2=40$ 但超出 $57$，应为... 右子编号 $2i+2=40$。
\end{solution}
\end{question}

\begin{question}
有 $n$ 个数据对象的二路归并排序中，每趟归并的时间复杂度为 \underline{\hspace{4em}}。
\begin{solution}
每趟归并需遍历所有元素一次，时间复杂度 $O(n)$。
\end{solution}
\end{question}

\begin{question}
对一组记录进行降序排序，其关键码为 $(46,70,56,38,40,80,60,22)$，采用初始步长为 $4$ 的希尔排序，第一趟扫描的结果是 \underline{\hspace{4em}}，而采用归并排序第一轮归并的结果是 \underline{\hspace{4em}}。
\begin{solution}
希尔排序步长 $4$：比较相距 $4$ 的元素并排序；归并排序第一轮两两归并。
\end{solution}
\end{question}

\begin{question}
如果只想得到 $1000$ 个元素的序列中最小的前 $5$ 个元素，在冒泡排序、快速排序、堆排序和归并排序中，哪种算法最快？\underline{\hspace{4em}}
\begin{solution}
用堆排序建堆后取 $5$ 次堆顶，时间复杂度 $O(n+5\log n)$。
\end{solution}
\end{question}

\begin{question}
如果一个图节点多而边少（稀疏图），适宜采用 \underline{\hspace{4em}} 方式进行存储。
\begin{solution}
邻接表空间复杂度 $O(n+e)$，邻接矩阵 $O(n^2)$，稀疏图用邻接表更省空间。
\end{solution}
\end{question}

\clearpage

\textbf{四、简答题（24 分，每小题 6 分）}

\begin{question}[森林周游]（6 分）

树周游算法可以很好地应用到森林的周游上。查看下列森林结构，请给出其深度优先周游序列和广度优先周游序列。

\begin{solution}
深度优先（先根遍历）：依次遍历每棵树的根结点，再递归遍历其子树。广度优先（层次遍历）：按层从上到下、从左到右遍历所有节点。具体序列需根据森林的实际结构确定。
\end{solution}
\end{question}

\begin{question}[Huffman树]（6 分）

设给定五个字符，其相应的权值分别为 $\{4,8,6,9,18\}$，试画出相应的哈夫曼树，并计算它的带权外部路径长度 WPL。

\begin{solution}
Huffman 树构建：每次合并权值最小的两个节点。
\begin{itemize}
  \item 合并 $4,6$ 得 $10$
  \item 合并 $8,9$ 得 $17$
  \item 合并 $10,17$ 得 $27$
  \item 合并 $18,27$ 得 $45$
\end{itemize}
WPL = $(4+6)\times3 + (8+9)\times3 + 18\times2 = 30+51+36=117$。
\end{solution}
\end{question}

\begin{question}[堆操作]（6 分）

下图是一棵完全二叉树：
\pyfig[0.24\linewidth]{figures/数据结构与算法B-21春-期末-fig01.png}{}

\begin{enumerate}[label=\arabic*)]
  \item 请根据初始建堆算法对该完全二叉树建堆，请画出构建的小根堆；（2 分）
  \item 基于（1）中得到的堆，删除其中的最小元素，请用图给出堆的调整过程；（2 分）
  \item 基于（1）中得到的堆，向其中插入元素 $2$，请给出堆的调整过程。（2 分）
\end{enumerate}

\begin{solution}
\begin{enumerate}[label=\arabic*)]
  \item 从最后一个非叶节点开始向下调整，构建小根堆。
  \item 删除堆顶（最小元素），将末尾元素移至堆顶，向下调整恢复堆性质。
  \item 在堆末尾插入 $2$，向上调整（与父节点比较交换）恢复堆性质。
\end{enumerate}
\end{solution}
\end{question}

\begin{question}[Prim算法求最小生成树]（6 分）

已知图 $G$ 的顶点集合 $V=\{V_0,V_1,V_2,V_3,V_4\}$，邻接矩阵如下（$\infty$ 表示无边）：
\[
\begin{pmatrix}
0 & 7 & \infty & 4 & 2\\
7 & 0 & 9 & 1 & 5\\
\infty & 9 & 0 & 3 & \infty\\
4 & 1 & 3 & 0 & 10\\
2 & 5 & \infty & 10 & 0
\end{pmatrix}
\]

\begin{enumerate}[label=\arabic*)]
  \item 根据 Prim 算法，求图 $G$ 从顶点 $V_0$ 出发的最小生成树；（2 分）
  \item 给出每一步执行后的 mst 数组。（4 分）
\end{enumerate}

\begin{solution}
从 $V_0$ 出发：
\begin{itemize}
  \item 选 $V_0-V_4$（权 $2$），加入 $V_4$
  \item 选 $V_0-V_3$（权 $4$）或 $V_4-V_1$（权 $5$），取最小 $V_0-V_3$（权 $4$），加入 $V_3$
  \item 选 $V_3-V_1$（权 $1$），加入 $V_1$
  \item 选 $V_3-V_2$（权 $3$），加入 $V_2$
\end{itemize}
最小生成树：$\{(V_0,V_4),(V_0,V_3),(V_3,V_1),(V_3,V_2)\}$，总权值 $2+4+1+3=10$。
\end{solution}
\end{question}

\clearpage

\textbf{五、算法题（16 分，每小题 8 分）}

\begin{question}[根据前序求后序]（8 分）

下列 $f$ 函数的功能是根据一个满二叉树的前序遍历序列，求其后序遍历序列，请完成填空（假设序列长度不超过 $32$）。

\begin{quote}
\footnotesize
\begin{tabular}{@{}l@{}}
\#define MX 40\\
void f(char *seq, int s, int L, char *result)\\
\{\\
    memset(result, 0, MX);\\
    char seq1[MX] = \{\}, seq2[MX] = \{\}, seq3[MX] = \{\};\\
    if (L == 1) \{\\
\underline{\hspace{12em}} /* (1) 1分 */ */\\
        return;\\
    \}\\
    int s1 = s + 1;\\
\underline{\hspace{12em}} /* (2) 2分 */ */\\
\underline{\hspace{12em}} /* (3) 1分 */ */\\
    f(seq, s1, L, seq1);                   /* 递归处理左子树 */\\
\underline{\hspace{12em}} /* (4) 2分 */ */\\
    strcpy(result, seq1);\\
\underline{\hspace{12em}} /* (5) 2分 */ */\\
    strcat(result, seq3);\\
\}\\
\end{tabular}
\end{quote}

\begin{solution}
填空答案：(1) \texttt{result[0] = seq[s]} (2) \texttt{L / 2} (3) \texttt{s + L + 1} (4) \texttt{f(seq, s2, L, seq3)} (5) \texttt{strcat(result, seq2)}

注意：\texttt{seq2} 在原代码中未赋值，应改为 \texttt{result[0]=seq[s]; strcat(result, seq1); strcat(result, seq3);}（后序：左-右-根）。
\end{solution}
\end{question}

\begin{question}[深度优先周游]（8 分）

阅读下列程序，完成图的深度优先周游算法。图采用邻接表表示。

\begin{quote}
\footnotesize
\begin{tabular}{@{}l@{}}
typedef struct EdgeNode \{\\
    int endvex;\\
    struct EdgeNode *nextedge;\\
\} EdgeNode;\\
typedef struct \{\\
    VexType vertex;\\
    struct EdgeNode *edgelist;\\
\} VexNode;\\
typedef struct \{\\
    VexNode vexs[1000];\\
    int n;\\
\} GraphList;\\
\\
void dFSInList(GraphList *pgraphlist, int visited[], int i)\\
\{\\
    PEdgeNode p;\\
    printf("node: \%c\textbackslash\{\}n", pgraphlist->vexs[i].vertex);\\
    visited[i] = 1;\\
\underline{\hspace{12em}} /* (1) 2分 */ */\\
\underline{\hspace{12em}} /* (2) 2分 */ */\\
        if (visited[p->endvex] == FALSE)\\
\underline{\hspace{12em}} /* (3) 2分 */ */\\
        p = p->nextedge;\\
    \}\\
\}\\
\\
void traverDFS(GraphList *pgraphlist)\\
\{\\
    int visited[MAXVEX];\\
    int n = pgraphlist->n;\\
    for (int i = 0; i < n; i++) visited[i] = FALSE;\\
    for (int i = 0; i < n; i++)\\
\underline{\hspace{12em}} /* (4) 2分 */ */\\
            dFSInList(pgraphlist, visited, i);\\
\}\\
\end{tabular}
\end{quote}

\begin{solution}
填空答案：(1) \texttt{pgraphlist->vexs[i].edgelist} (2) \texttt{p != NULL} (3) \texttt{pgraphlist, visited, p->endvex} (4) \texttt{visited[i] == FALSE}
\end{solution}
\end{question}

\clearpage
